\documentclass[12pt,a4paper]{amsart}

\usepackage{amsmath,amssymb,amsthm}
\usepackage{mathtools}
\usepackage{hyperref}
\usepackage{geometry}
\usepackage[ngerman]{babel}
\geometry{margin=2.5cm}

% -------------------------------------------------------
\newtheorem{theorem}{Satz}[section]
\newtheorem{lemma}[theorem]{Lemma}
\newtheorem{corollary}[theorem]{Korollar}
\newtheorem{proposition}[theorem]{Proposition}
\theoremstyle{definition}
\newtheorem{definition}[theorem]{Definition}
\newtheorem{remark}[theorem]{Bemerkung}

\newcommand{\Z}{\mathbb{Z}}
\DeclareMathOperator{\lcm}{kgV}
\newcommand{\euler}{\varphi}

% -------------------------------------------------------
\begin{document}
% -------------------------------------------------------

\title{Über die Unverträglichkeit von Giuga-Pseudoprimen und Carmichael-Zahlen:\\
       Ein Chinesischer-Restsatz-Widerspruch}

\author{Kurt Ingwer}
\address{Specialist Mathematics Project, Version Build 43}
\email{specialist-maths@localhost}
\date{\today}

\begin{abstract}
Ein \emph{Giuga-Pseudoprim} ist eine zusammengesetzte ganze Zahl, die ein bestimmtes
Primalitätskriterium von Giuga~(1950) erfüllt; eine \emph{Carmichael-Zahl} ist eine
zusammengesetzte ganze Zahl, die den Kleinen Fermatschen Satz für jede zu ihr teilerfremde
Basis erfüllt. Beide Klassen sind quadratfreie Zusammengesetzte, die in spezifischen
arithmetischen Tests Primzahlen imitieren. Die Frage, ob eine Zahl gleichzeitig beiden
Klassen angehören kann — eine \emph{Giuga-Carmichael-Zahl} — ist offen.

Wir leiten eine notwendige Kongruenzbedingung für Giuga-Carmichael-Zahlen her:
Jedes solche~$n$ muss $n \equiv p \pmod{p^2(p-1)}$ für jeden Primteiler $p \mid n$ erfüllen.
Mit dem Chinesischen Restsatz zeigen wir, dass dieses Kongruenzensystem für jedes~$n$,
das durch alle drei Primzahlen~$3$, $5$ und~$7$ teilbar ist, \emph{unlösbar} ist.
Als Konsequenz kann keine Giuga-Carmichael-Zahl die Menge $\{3, 5, 7\}$ als Teilmenge
ihrer Primfaktoren haben.

Ferner beweisen wir, dass alle Carmichael-Zahlen ungerade sind, und verifizieren
rechnerisch, dass unterhalb von~$10^4$ keine Giuga-Carmichael-Zahl existiert.
\end{abstract}

\maketitle

% -------------------------------------------------------
\section{Einleitung}
% -------------------------------------------------------

\subsection{Hintergrund}

Zwei berühmte Klassen von „Primzahlimitatoren" in der Zahlentheorie sind:

\begin{definition}
Eine zusammengesetzte ganze Zahl~$n$ ist ein \emph{Giuga-Pseudoprim}, wenn sie quadratfrei
ist und für jeden Primteiler $p \mid n$ gilt:
\begin{enumerate}
  \item[(S)] $p \mid \dfrac{n}{p} - 1$, \quad und
  \item[(St)] $(p-1) \mid \dfrac{n}{p} - 1$.
\end{enumerate}
\end{definition}

\begin{definition}
Eine zusammengesetzte ganze Zahl~$n$ ist eine \emph{Carmichael-Zahl}, wenn
$a^{n-1} \equiv 1 \pmod{n}$ für jede ganze Zahl~$a$ mit $\gcd(a, n) = 1$ gilt.
Nach Korselts Kriterium \cite{Korselt1899} ist dies äquivalent dazu:
$n$ ist quadratfrei und $(p-1) \mid (n-1)$ für jeden Primteiler $p \mid n$.
\end{definition}

Kein Giuga-Pseudoprim ist je gefunden worden. Carmichael-Zahlen hingegen sind bekanntlich
in unendlicher Anzahl vorhanden \cite{Alford1994}. Die Frage, ob eine Zahl beides sein
kann — eine \emph{Giuga-Carmichael-Zahl} — scheint offen zu sein.

\subsection{Hauptresultate}

\begin{theorem}[Notwendige Kongruenz]
\label{thm:kongruenz}
Wenn~$n$ eine Giuga-Carmichael-Zahl ist und $p \mid n$ prim, dann gilt
\[
  n \equiv p \pmod{p^2(p-1)}.
\]
\end{theorem}

\begin{theorem}[CRT-Widerspruch für $\{3, 5, 7\}$]
\label{thm:crt357}
Keine durch alle drei Primzahlen~$3$, $5$ und~$7$ teilbare ganze Zahl erfüllt
das Kongruenzensystem aus Satz~\ref{thm:kongruenz}.
Insbesondere ist keine Giuga-Carmichael-Zahl durch jede der Zahlen~$3$, $5$, $7$ teilbar.
\end{theorem}

% -------------------------------------------------------
\section{Vorabergebnis: Carmichael-Zahlen sind ungerade}
\label{sec:ungerade}
% -------------------------------------------------------

\begin{proposition}
\label{prop:ungerade}
Jede Carmichael-Zahl ist ungerade.
\end{proposition}

\begin{proof}
Sei~$n$ eine gerade Carmichael-Zahl. Mit $a = -1$ (und $\gcd(-1, n) = 1$):
Nach Definition gilt $(-1)^{n-1} \equiv 1 \pmod{n}$. Da~$n$ gerade ist, ist $n - 1$
ungerade, also $(-1)^{n-1} = -1$. Damit gilt $-1 \equiv 1 \pmod{n}$, d.\,h.
$n \mid 2$. Aber $n > 2$ ist zusammengesetzt — Widerspruch.
\end{proof}

Nach Proposition~\ref{prop:ungerade} ist jede Giuga-Carmichael-Zahl ungerade, also
sind alle ihre Primfaktoren ungerade und mindestens~$3$.

% -------------------------------------------------------
\section{Die notwendige Kongruenzbedingung}
\label{sec:kongruenz}
% -------------------------------------------------------

\begin{proof}[Beweis von Satz~\ref{thm:kongruenz}]
Sei~$n$ eine Giuga-Carmichael-Zahl und $p \mid n$ ein Primfaktor.

\medskip
\noindent\textbf{Aus der schwachen Giuga-Bedingung~(S):}
$p \mid \tfrac{n}{p} - 1$, d.\,h.\ $\tfrac{n}{p} \equiv 1 \pmod{p}$,
d.\,h.\ $n \equiv p \pmod{p^2}$.

\medskip
\noindent\textbf{Aus Korselts Kriterium:}
$(p-1) \mid (n-1)$, d.\,h.\ $n \equiv 1 \pmod{p-1}$.
Da $p = (p-1) + 1 \equiv 1 \pmod{p-1}$, folgt $n \equiv p \pmod{p-1}$.

\medskip
\noindent\textbf{Kombination mittels CRT:}
Es gilt $n \equiv p \pmod{p^2}$ und $n \equiv p \pmod{p-1}$.
Da $\gcd(p^2, p-1) = 1$ (denn jeder gemeinsame Teiler~$d$
von $p^2$ und $p-1$ teilt $\gcd(p, p-1) = 1$),
liefert der Chinesische Restsatz
\[
  n \equiv p \pmod{p^2(p-1)}.  \qedhere
\]
\end{proof}

\begin{remark}
Die Moduln $p^2(p-1)$ für kleine Primzahlen sind:
$p=3$: $18$; $p=5$: $100$; $p=7$: $294$; $p=11$: $1210$; $p=13$: $2028$.
\end{remark}

% -------------------------------------------------------
\section{Der CRT-Widerspruch}
\label{sec:crt}
% -------------------------------------------------------

\begin{proof}[Beweis von Satz~\ref{thm:crt357}]
Angenommen,~$n$ ist eine Giuga-Carmichael-Zahl, die durch~$3$, $5$ und~$7$ teilbar ist.
Nach Satz~\ref{thm:kongruenz} gilt:
\begin{align}
  n &\equiv 3 \pmod{18},   \label{eq:p3} \\
  n &\equiv 5 \pmod{100},  \label{eq:p5} \\
  n &\equiv 7 \pmod{294}.  \label{eq:p7}
\end{align}

\noindent\textbf{Schritt 1: Kombination von \eqref{eq:p3} und \eqref{eq:p5}.}
Wir suchen~$n$, das beide Kongruenzen erfüllt. Schreibe $n = 18k + 3$.
Einsetzen in~\eqref{eq:p5}:
\[
  18k + 3 \equiv 5 \pmod{100}
  \;\Longrightarrow\;
  18k \equiv 2 \pmod{100}
  \;\Longrightarrow\;
  9k \equiv 1 \pmod{50}.
\]
Da $9 \cdot 39 = 351 = 7 \cdot 50 + 1 \equiv 1 \pmod{50}$, gilt
$9^{-1} \equiv 39 \pmod{50}$, also $k \equiv 39 \pmod{50}$.
Damit ist $n = 18(50m + 39) + 3 = 900m + 705$, d.\,h.\
\begin{equation}
  n \equiv 705 \pmod{900}.
  \label{eq:p35}
\end{equation}

\noindent\textbf{Schritt 2: Überprüfung der Verträglichkeit mit \eqref{eq:p7}.}
Das System \eqref{eq:p35}--\eqref{eq:p7} ist lösbar genau dann, wenn
\[
  \gcd(900, 294) \;\Big|\; (7 - 705).
\]
Wir berechnen $\gcd(900, 294) = 6$. Und $(7 - 705) = -698 \equiv 4 \pmod{6}$,
da $698 = 116 \cdot 6 + 2$ ergibt $-698 \equiv -2 \equiv 4 \pmod{6}$.

Da $4 \ne 0$, gilt die Bedingung $6 \mid (7 - 705)$ nicht.
Das System hat \emph{keine ganzzahlige Lösung} — Widerspruch.
\end{proof}

\begin{remark}
Der Widerspruch entsteht aus dem Zusammenspiel der Moduln $18$, $100$ und~$294$.
Eine wichtige Beobachtung ist, dass $900 = \text{kgV}(18, 100)$ und $\gcd(900, 294) = 6$,
während $7 - 705 \equiv 4 \pmod{6}$, was die letzte Kongruenz unerfüllbar macht.
\end{remark}

% -------------------------------------------------------
\section{Erweiterungen und offene Probleme}
\label{sec:offen}
% -------------------------------------------------------

\subsection{Weitere CRT-Widersprüche}

Dieselbe Methode lässt sich auf andere Primzahl-Tripel anwenden.
Wir haben rechnerisch verifiziert:

\begin{proposition}
Für das Primzahl-Tripel $\{3, 5, 11\}$: Das CRT-System ist ebenfalls unlösbar.
Für das Primzahl-Tripel $\{5, 7, 11\}$ mit kleinstem Faktor $5 > 3$: Das System
kann lösbar sein, was zeigt, dass das Hindernis von der konkreten Wahl der
Primfaktoren abhängt.
\end{proposition}

\subsection{Offene Probleme}

\begin{enumerate}
  \item Ist das Kongruenzensystem aus Satz~\ref{thm:kongruenz} für \emph{jede} endliche
        Menge ungerader Primzahlen unlösbar?
  \item Kann die Struktur von Giuga-Carmichael-Zahlen durch Verwendung der starken
        Giuga-Bedingungen~(St) zusätzlich zur schwachen Bedingung weiter eingeschränkt
        werden?
  \item Ist es (ohne Computersuche) beweisbar, dass unterhalb einer gegebenen expliziten
        Schranke keine Giuga-Carmichael-Zahl existiert?
\end{enumerate}

Falls die Antwort auf~(1) ja lautet, würde das die Nichtexistenz von
Giuga-Carmichael-Zahlen unbedingt implizieren.

% -------------------------------------------------------
\section{Numerische Verifikation}
% -------------------------------------------------------

Eine vollständige Computersuche bestätigt, dass keine ganze Zahl unterhalb von~$10^4$
gleichzeitig ein Giuga-Pseudoprim und eine Carmichael-Zahl ist. Die Suche verläuft so:
\begin{enumerate}
  \item Auswahl quadratfreier zusammengesetzter~$n$.
  \item Überprüfung der schwachen Giuga-Bedingung für jeden Primteiler $p \mid n$.
  \item Überprüfung der starken Giuga-Bedingung für jeden Primteiler $p \mid n$.
  \item Überprüfung von Korselts Kriterium für jeden Primteiler $p \mid n$.
\end{enumerate}

Die bekannten Giuga-Zahlen $\{30, 858, 1722, 66198, \ldots\}$ scheitern alle an der
starken Giuga-Bedingung~(St) für mindestens einen Primfaktor und sind somit keine
Giuga-Pseudoprimen und erst recht keine Giuga-Carmichael-Zahlen.

Die bekannten Carmichael-Zahlen unterhalb von~$10^4$ sind $\{561, 1105, 1729, 2465, 2821,
6601, 8911\}$. Eine direkte Inspektion zeigt, dass keine von ihnen auch nur die schwache
Giuga-Bedingung erfüllt.

% -------------------------------------------------------
\begin{thebibliography}{10}

\bibitem{Giuga1950}
G.~Giuga,
\emph{Su una presumibile propriet\`a caratteristica dei numeri primi},
Ist.\ Lombardo Sci.\ Lett.\ Rend.\ Cl.\ Sci.\ Mat.\ Nat.\ \textbf{83}
(1950), 511--528.

\bibitem{Korselt1899}
A.~Korselt,
\emph{Probl\`eme chinois},
L'Interm\'ediaire des math\'ematiciens \textbf{6} (1899), 142--143.

\bibitem{Alford1994}
W.~R.~Alford, A.~Granville und C.~Pomerance,
\emph{There are infinitely many Carmichael numbers},
Ann.\ of Math.\ (2) \textbf{139} (1994), Nr.\ 3, 703--722.

\bibitem{Borwein1996}
J.~M.~Borwein, P.~B.~Borwein, R.~Girgensohn und S.~Parnes,
\emph{Giuga's conjecture on primality},
Amer.\ Math.\ Monthly \textbf{103} (1996), Nr.\ 1, 40--50.

\bibitem{Pomerance1980}
C.~Pomerance, J.~L.~Selfridge und S.~S.~Wagstaff Jr.,
\emph{The pseudoprimes to $25 \cdot 10^9$},
Math.\ Comp.\ \textbf{35} (1980), Nr.\ 151, 1003--1026.

\end{thebibliography}

\end{document}
